\documentclass[10pt,a4paper]{amsart}
\usepackage[margin=19mm]{geometry}
\usepackage{amsmath,amssymb,mathtools}
\usepackage[scr=rsfs,cal=euler]{mathalfa}
\usepackage{xfrac}
\usepackage[unicode,hidelinks,bookmarks=false]{hyperref}
\newcommand{\ie}{i.e.\ }
\newcommand{\cf}{cf.\ }
\newcommand{\resp}{respectively\ }
\newcommand{\Subsection}{\subsection}
\providecommand{\nobreakdash}{\nobreak\hbox{-}}
\setlength{\emergencystretch}{2em}
\multlinegap0pt
\usepackage{bm}
\usepackage{algorithm}\usepackage{algpseudocode}
\usepackage{tikz}
\usepackage{xcolor}
\usepackage{amssymb}
\usepackage{dsfont}
\usepackage{enumerate}
\newtheorem{prop}{Proposition}
\newtheorem{lemma}{Lemma}
	\newcommand{\cX}{\mathcal{X}}
\title{Optimal Multi-bit Generative Watermarking Schemes Under Worst-Case False-Alarm Constraints}
\author{Yu-Shin Huang, Chao Tian, Krishna Narayanan}
\date{Source: arXiv:2604.08759v1 (CC BY 4.0)}
\begin{document}
\maketitle

\noindent\emph{Source and licence.} Optimal Multi-bit Generative Watermarking Schemes Under Worst-Case False-Alarm Constraints. Version arXiv:2604.08759v1 (CC BY 4.0), distributed by arXiv under the Creative Commons Attribution 4.0 International licence, http://creativecommons.org/licenses/by/4.0/, as recorded in the arXiv licence metadata for this exact version and shown on the abstract page https://arxiv.org/abs/2604.08759v1. Full licence notice: \texttt{licenses/arxiv-2604.08759-CC-BY-4.0.txt}.\par\smallskip

\noindent\textit{Editorial scope.} Counted unit: the article's lemma lemma:extend\_PX (source0.tex lines 956--958). The article's complete original proof of it is delivered with it (source0.tex lines 959--982, one \texttt{proof} environment of 24 lines). No mathematics is written, repaired or re-derived: the statement and the proof are the article's own lines, in the article's own notation, and no retained line is substituted or re-flowed.  Retained uncounted as the source-local context the proof needs: lines 526--528; lines 568--579; lines 876--888.  Nothing was omitted from any retained span, so the package carries no omission record.  No retained line contains a citation, so the wrapper carries no bibliography and no bibliography entry.  No retained line contains a figure environment, an image-inclusion command or drawing code, so no image is delivered and the package declares no asset dependency.  Editorial formatting only, in addition: an AMS \texttt{amsart} preamble that restates the article's own declarations and macros used by the retained lines, namely the environments \texttt{prop}, \texttt{lemma} on separate counters, together with the standard packages those lines need.  The retained environments print in this document as Proposition 1, Lemma 1 in order of appearance, whereas the article prints the same items with the numbering of its own full document.  The complete native source of the article as distributed in the arXiv e-print for this exact version is bundled unchanged as \texttt{sources/198169/source0.tex}.
    \begin{align}
        T \cdot \max_i \boldsymbol{a}(i) \leq \sum_{i = 1}^N \boldsymbol{a}(i).\label{eq:T-hot-representable}
    \end{align}

\begin{prop}[$P_X$ Decomposition] \label{prop:split_px}
    Let $P_X$ be a distribution with $P_X(1) \leq P_X(2) \leq \dots \leq P_X(N)$. There exists a decomposition $P_X^{(1)}+P_X^{(2)}+P_X^{(3)} = P_X$ such that 
    \begin{itemize}
        \item[1.] $P_X^{(1)}$ is $T$-hot representable;
        \item[2.] $P_X^{(2)}$ is a non-negative non-decreasing vector with at most $T-1$  positive components, all of which are located in the final entries;
        \item[3.] $P_X^{(3)}(x) = \max(P_X(x) -\tfrac{\alpha}{T}, 0), \forall x$.
    \end{itemize}
    Moreover, either $P^{(2)}_X(x)=0$ for all $x\in\cX$, or 
    \begin{align}
    T\cdot\max_{x\in \cX}P_X^{(1)}(x)=\sum_{x\in\cX}P_X^{(1)}(x).\label{eqn:PX2}
    \end{align}
\end{prop}

\begin{algorithm}[h]
\caption{$P_m^{(3)}$ Construction}
\label{alg:pm3_proportional_fill}
\begin{algorithmic}[1]
\State \textbf{Input:} $P_X^{(3)}$ with $\tilde{K}$ positive value; Decomposition $P_X^{(2)} = \sum_{j=1}^K \Delta\eta_{N-j+1}\boldsymbol{u}_j$; $K$ number of nonzero elements in $P_X^{(2)}$; Total Imbalance $U$; Initialized $P_m^{(3)} =\mathbf{0}$, $m \in [1:T]$
\For{$j = K, \cdots, 1$}
    \For{$x = N-\tilde{K}+1, \ldots, N$}
        \State $P_m^{(3)}(x,\zeta) \gets P_m^{(3)}(x,\zeta) + \frac{\Delta\eta_{N-j+1}(T-j)}{|S(K) \setminus \cup_{l=N-j+1}^N \gamma^{-1}_l(m)|}\times \frac{P_X^{(3)}(x)}{R}$ , if $\zeta \in S(K) \setminus \cup_{l=N-j+1}^N \gamma^{-1}_l(m)  $ \label{alg-eq:allocate_imbalance}
    \EndFor
\EndFor
\State $P_m^{(3)}(x, \mathbf{0}_N) \gets P_X^{(3)}(x) \times (1-\frac{U}{R}), \forall x$
\end{algorithmic}
\end{algorithm}

\begin{lemma} \label{lemma:extend_PX}
    Let \(P_X\) be a non-decreasing distribution of length \(N\). Then there exists an integer \(n \ge 0\) such that we can construct an extended non-negative vector \(P'_X\) of length \(N + n\) that is \(T\)-hot representable. Specifically, when \(n = 0\), we have \(P'_X = \min \bigl(\tfrac{\alpha}{T}, P_X\bigr)\). When \(n > 0\), the components of \(P'_X\) add up to \(1\), i.e. \(P'_X\) itself forms a probability distribution with extra $n$ psuedo-tokens. 
\end{lemma}

\begin{proof}
    We define $\boldsymbol{a}$ and $\boldsymbol{r}$ as
    \begin{align}
        \boldsymbol{a}(x) = \min\!\left(\frac{\alpha}{T}, P_X(x)\right), \; \boldsymbol{r}(x) = (P_X - \boldsymbol{a})(x) = \max\!\left(P_X(x) - \frac{\alpha}{T}, 0\right),
    \end{align}
    where $\boldsymbol{a}$ is same as what we defined in the proof of Proposition~\ref{prop:split_px}, and $\boldsymbol{r}$ is the same definition as $P_X^{(3)}$. We use $R$ to denote the sum of the vector $\boldsymbol{r}$, consistent with the definition of $R$ provided in Algorithm \ref{alg:pm3_proportional_fill}.
    
    If $\boldsymbol{a}$ itself is $T$-hot representable, then set $n = 0$ and $P'_X = \boldsymbol{a}$. Otherwise, we choose $n = \lceil \frac{R}{\boldsymbol{a}(N)} \rceil$. We then demonstrate that
    \begin{align}
        P'_X = (\boldsymbol{a}, \tfrac{R}{n}\mathbf{1}_n)
    \end{align}
    is $T$-hot representable, where $\mathbf{1}_n$ denotes the all-ones vector of length $n$. The sum of $P_X'$ is 
    \begin{align*}
        \sum_{i=1}^{N+n} P_X'(i) & = \sum_{i=1}^{N} \boldsymbol{a}(i) + \sum_{i=N+1}^{N+n} \frac{R}{n} = \sum_{i=1}^{N} \boldsymbol{a}(i) + R = \sum_{i=1}^{N} \boldsymbol{a}(i) + \sum_{i=1}^{N}[P_X(i) - \boldsymbol{a}(i)]  = 1.
    \end{align*}
    
    Clearly, we have $\tfrac{R}{n} \leq \tfrac{R}{\tfrac{R}{\boldsymbol{a}(N)}} = \boldsymbol{a}(N)$. Furthermore, because $P_X$ is a non-decreasing vector, the vector $\boldsymbol{a}$ is also non-decreasing. Hence, $\max_i P_X'(i) = \boldsymbol{a}(N)$.

    The vector $P_X'$ is $T$-hot representable since
    \begin{align}
        T \max_i P_X'(i) = T\boldsymbol{a}(N) \leq T \frac{\alpha}{T} = \alpha \leq 1 = \sum_{i=1}^{N+n} P_X'(i),
    \end{align}
    which satisfies the inequality \eqref{eq:T-hot-representable}.
\end{proof}

\end{document}
